\let\negmedspace\undefined
\let\negthickspace\undefined
\documentclass[journal]{IEEEtran}
\usepackage[a5paper, margin=10mm, onecolumn]{geometry}
%\usepackage{lmodern} % Ensure lmodern is loaded for pdflatex
\usepackage{tfrupee} % Include tfrupee package

\setlength{\headheight}{1cm} % Set the height of the header box
\setlength{\headsep}{0mm}     % Set the distance between the header box and the top of the text

\usepackage{gvv-book}
\usepackage{gvv}
\usepackage{cite}
\usepackage{amsmath,amssymb,amsfonts,amsthm}
\usepackage{algorithmic}
\usepackage{graphicx}
\usepackage{textcomp}
\usepackage{xcolor}
\usepackage{txfonts}
\usepackage{listings}
\usepackage{enumitem}
\usepackage{mathtools}
\usepackage{gensymb}
\usepackage{comment}
\usepackage[breaklinks=true]{hyperref}
\usepackage{tkz-euclide} 
\usepackage{listings}
% \usepackage{gvv}                                        
\def\inputGnumericTable{}                                 
\usepackage[latin1]{inputenc}                                
\usepackage{color}                                            
\usepackage{array}                                            
\usepackage{longtable}                                       
\usepackage{calc}                                             
\usepackage{multirow}                                         
\usepackage{hhline}                                           
\usepackage{ifthen}                                           
\usepackage{lscape}
\begin{document}

\bibliographystyle{IEEEtran}
\vspace{3cm}

\title{2015-XE-14-26}
\author{EE24BTECH11065 - spoorthi yellamanchali}
% \maketitle
% \newpage
% \bigskip
{\let\newpage\relax\maketitle}

\renewcommand{\thefigure}{\theenumi}
\renewcommand{\thetable}{\theenumi}
\setlength{\intextsep}{10pt} % Space between text and floats


\numberwithin{equation}{enumi}
\numberwithin{figure}{enumi}
\renewcommand{\thetable}{\theenumi}


\begin{enumerate}
    \item For a complex number $k$ and a complex variable $z$, the complex function $(z + k)^2$ is analytic \hfill{[XE-2015]}
    \begin{enumerate}
        \item for all $k$.
        \item only for all $k$ whose imaginary component is zero.
        \item only for all $k$ whose real component is zero.
        \item only when $k$ is zero.
    \end{enumerate}
\item The divergence of a vector field $\vec{V}(x, y, z) = 2^x \, \hat{i} + 2^y \, \hat{j} + 2^z \, \hat{k}$ at a point $(1,1,1)$ is \\.\hfill{[XE-2015]} \\
    \item The type of the differential equation $(1 - x) \frac{d^3 y}{dx^3} = \sqrt{1 + \left( \frac{dy}{dx} \right)^2} + 5y = \cos(x)$ is \\. \hfill{[XE-2015]}
    \begin{enumerate}
        \item linear and first order.
        \item non-linear and first order.
        \item linear and third order.
        \item non-linear and third order.
    \end{enumerate}

    \item A box has ten light bulbs out of which two are defective. Two light bulbs are drawn from this box one after the other without replacement. The probability that both light bulbs drawn are NOT defective is \hfill{[XE-2015]}
    \begin{multicols}{4}
    \begin{enumerate}
        \item $\frac{8}{45}$
        \item $\frac{28}{45}$
        \item $\frac{16}{25}$
        \item $\frac{4}{5}$
    \end{enumerate}
    \end{multicols}

   \item The value of the following limit is 
\[
\lim_{x \to \infty} \left( \frac{x - \sin(x)}{x + \sin(x)} + \frac{\ln x}{x} \right)
\] \hfill{[XE-2015]}
\begin{multicols}{4}
    

\begin{enumerate}
    \item zero
    \item 1
    \item 2
    \item $\infty$
\end{enumerate}
\end{multicols}
    \item The value of the Fourier coefficient, $A_0$, in the series 
    
    \[  \left\{ A_0 + \sum_{n=1}^{\infty} \left( a_n \cos(nx) + b_n \sin(nx) \right) \right\}
   \]
    of a function $f(x) = x^2$ with a period $2\pi$ defined over an interval $0 \leq x \leq 2\pi$ is \\. \hfill{[XE-2015]}
    \begin{multicols}{4}
        
    
    \begin{enumerate}
        \item $\frac{4 \pi^2}{3}$
        \item $\frac{2 \pi^2}{3}$
        \item $\frac{\pi^2}{3}$
        \item $\frac{\pi^2}{6}$
    \end{enumerate}
    \end{multicols}

    \item The general solution, $y(x)$, for the differential equation 
\[
    x \frac{d^2 y}{dx^2} - \frac{dy}{dx} - 1 = 0
\]
    ($c_1$ and $c_2$ are real constants) is \hfill{[XE-2015]}
    \begin{multicols}{2}
        
    
    \begin{enumerate}
        \item $c_1 \frac{x^2}{2} + 2x + c_2$
        \item $c_1 \frac{x^2}{2} - x + c_2$
        \item $c_1 \frac{x^2}{2} + x + c_2$
        \item $c_1 \frac{x^2}{2} - 2x + c_2$
    \end{enumerate}
\end{multicols}
    \item The minimum number of iterations required to evaluate $\sqrt{28}$, correct up to three decimal places, by the Newton-Raphson method with an initial guess of 5 is \\.\hfill{[XE-2015]} \\
    


        \item  The gap $\delta$ between two concentric cylinders, each of height $h$, is filled with an oil. The torque required to rotate the inner cylinder at an angular velocity of $\omega$ against the fixed outer cylinder is $T$. The diameter of the inner cylinder is $d$ and $\delta \ll d$. The dynamic viscosity of the oil is given by \hfill{[XE-2015]}
    \begin{multicols}{4}
        
    
    \begin{enumerate}
        \item $\frac{4 \pi \delta T}{d^3 \omega h}$
        \item $\frac{4 \delta T}{\pi d^3 \omega h^2}$
        \item $\frac{4 \pi \delta T}{d^2 \omega h^2}$
        \item $\frac{4 \delta T}{\pi d \omega h^3}$
    \end{enumerate}
    \end{multicols}
    \item Water is retained against a sluice gate in the form of a circular segment as shown in the figure.If $\rho$ and $g$ are the density of water and gravitational acceleration respectively, the upward force exerted by the gate on the water per unit depth perpendicular to the plane of the figure is \hfill{[XE-2015]} \\
    
   
\begin{figure}[H]
\centering
\resizebox{0.5\textwidth}{!}{%
\begin{circuitikz}[scale =0.5]
\tikzstyle{every node}=[font=\small]
\draw [dashed] (15.5,22) -- (15.5,15);
\draw [dashed] (15.5,22) -- (19.25,18.25);
\draw [dashed] (19.25,18.25) -- (15.5,14.5);
\draw [dashed] (15.5,15) -- (15.5,14.25);
\draw [dashed] (19.25,18.25) -- (11.75,18.25);
\draw(15.5,22) arc[start angle =90,end angle =180,radius = 3.75];
\draw(11.75,18.25) arc[start angle=180,end angle = 270,radius = 3.75];
\draw(18,18.75) node {$\theta$};
\draw(18,17.75) node{$\theta$};
\draw [short] (15.5,22) -- (9.25,22);
\draw [short] (15.5,14.5) -- (9.25,14.5);
\draw [short] (10.75,21.75) -- (9.25,21.75);
\draw [short] (10.5,21.5) -- (9.5,21.5);
\draw [short] (10.25,21.25) -- (9.75,21.25);
\draw [short] (9.25,22) -- (8.75,22);
\draw [short] (10,22) -- (10.75,22.5);
\draw [short] (10.75,22.5) -- (9.25,22.5);
\draw [short] (9.25,22.5) -- (10,22);
\draw [short] (15.5,14.5) -- (16,14.5);
\draw [short] (15.25,14.5) -- (14.75,13.5);
\draw [short] (16,14.5) -- (16.5,13.5);
\draw [short] (14.75,13.5) -- (16.5,13.5);
\draw [<->, >=Stealth] (11.5,22) -- (11.5,14.5);
\draw [<->, >=Stealth] (18.5,19) -- (18,18.25);
\draw [<->, >=Stealth] (18,18.25) -- (18.5,17.5);
\draw [line width=0.7pt, short] (15.25,19.5) -- (18.75,21.5);
\draw [line width=0.7pt, ->, >=Stealth] (15.25,19.5) -- (12,19.5);
\node [font=\small ]at (10,18.5) {Water};
\node [font=\small] at (20,21.5) {circular };
\node [font=\small] at (19.75,21) {segment};
\node [font=\small]at (19.75,20.5) {of radius R};
\end{circuitikz}
}%

\label{fig:my_label}
\end{figure}
    \begin{multicols}{2}
    \begin{enumerate}
        
    \item $\rho R^2 \left( \theta - \frac{1}{2} \sin 2 \theta \right) g$
    \item $\rho R^2 \left( \cos^2 \theta - \frac{1}{2} \sin \theta \right) g$
    \item $\rho R^2 \left( \cos \theta - \frac{1}{2} \sin \theta \right) g$
    \item $\rho R^2 \left( \cos^2 \theta - \frac{1}{2} \sin^2 \theta \right) g$
\end{enumerate}
        
    
        
    \end{multicols}

   
\item A two-dimensional velocity field is given by $\vec{V} = 10(y^3 - x^2 y)\hat{i} + 2Cxy^2\hat{j}$, where $\hat{i}$ and $\hat{j}$ are the unit vectors in the directions of the rectangular Cartesian coordinates $x$ and $y$, respectively. If the flow is incompressible, the constant $C$ should be \hfill{[XE-2015]}
\begin{multicols}{4}
    

\begin{enumerate}
    \item $-10$
    \item $0$
    \item $5$
    \item $10$
\end{enumerate}
\end{multicols}

\item Let $\vec{V}$ and $T$ denote the velocity vector and temperature in a flow field. The rate of change of temperature experienced by a fluid particle as it passes through the point $(x_1, y_1, z_1)$ at a time $t_1$ is $2.5 \ ^\circ C/s$. The rate of change of temperature at a time $t_1$ at the point $(x_1, y_1, z_1)$ is $4.8 \ ^\circ C/s$. The quantity $\vec{V} \cdot \nabla T$ at $(x_1, y_1, z_1)$ and $t_1$ in $^\circ C/s$ is \underline{\hspace{2cm}}. \hfill{[XE-2015]}\\



\item In a simple Couette flow apparatus, the gap $h$ between the parallel plates is filled with a liquid of density $\rho$ and dynamic viscosity $\mu$, and one plate is dragged at a velocity of $U$ parallel to itself, while the other plate is fixed. The magnitude of vorticity at any point in the flow is \hfill{[XE-2015]} 
\begin{multicols}{4}
    

\begin{enumerate}
    \item $\frac{\mu}{\rho h^2}$
    \item $0$
    \item $\frac{1}{h^2} \sqrt{\frac{\mu U h}{\rho}}$
    \item $\frac{U}{h}$
\end{enumerate}

\end{multicols}

\end{enumerate}
  

\end{document}



